\chapter{Backtest-Engine Architecture}\label{ch:pl:backtest-engine-architecture}

A quoting strategy looked rich in the bar backtest, flat in the order-book replay, and lost money in the simulator. The firm had
four backtesters, written by four teams in four styles, each with its own idea of a strategy, a clock, an order and a result; the
strategy had been ported to each by hand. It took a week to be sure that the three answers differed because the markets differed
--- bars against queues, queues against a market that reacts --- and not because the four ports of the strategy did. This chapter
builds the engine that makes that question a query: one strategy object, one \omterm{def:pl:backtest-engine-architecture:event}{event model}, one \omterm{def:pl:backtest-engine-architecture:dal}{data-access layer} and one \omterm{def:pl:backtest-engine-architecture:results}{results
store} over the four fidelity levels of One Quant Book~7 (chapter~16), the fourth being a \omterm{def:pl:backtest-engine-architecture:reactive}{reactive simulation} on the exchange simulator
of One Quant Book~10.

\section{One engine, four levels}

\begin{definition}[Backtest engine]\label{def:pl:backtest-engine-architecture:engine}
A \emph{backtest engine}\index{backtest engine} is the platform component that runs a strategy against historical or simulated
markets at a chosen fidelity level through one strategy interface, reads its data through a \omterm{def:pl:backtest-engine-architecture:dal}{data-access layer}, records what the
strategy saw and did in one \omterm{def:pl:backtest-engine-architecture:event}{event model}, and stores every run with the identity of its code, data, parameters and seed.
\end{definition}

The firm's backtesters already exist: Book~7's vectorised backtester (\texttt{firm.vecbt}, level~1), its event-driven backtester on
bars (\texttt{firm.evbt}, level~2) and its order-book replay with queue positions (\texttt{firm.lobreplay}, level~3), and Book~10's
exchange simulator (\texttt{firm.exchsim}). Book~7 defined level~4 as live trading at small size or paper trading, and in its
chapter~19 let an agent trade inside the synthetic market as the stand-in for live; here the exchange simulator plays that part, with
the same session replayed as real orders against which the strategy's own orders rest and trade. The engine does not rewrite any of
them. Each level has an adapter that turns the one strategy object into what that backtester expects, and turns what the backtester
reports into the engine's events (\cref{fig:pl:backtest-engine-architecture:arch}).

\begin{omfigure}
\begin{tikzpicture}[font=\scriptsize,box/.style={draw,rounded corners=2pt,align=center,minimum height=0.75cm,inner sep=3pt},
  >=Stealth]
\node[box,fill=omMeth!15,minimum width=3.2cm] (s) at (0,2.6) {strategy object\\ \texttt{QuoteStrategy} or \texttt{TargetStrategy}};
\node[box,fill=omDef!15,minimum width=10.6cm] (e) at (2.6,1.2) {\texttt{Engine(level, strategy, data, config).run()}: one context
(\texttt{send}, \texttt{cancel}, \texttt{working}, \texttt{position}), one event log};
\foreach \x/\l/\c in {-2.1/{level 1\\ \texttt{firm.vecbt}}/1,0.9/{level 2\\ \texttt{firm.evbt}}/2,3.9/{level 3\\
  \texttt{firm.lobreplay}}/3,6.9/{level 4\\ \texttt{firm.exchsim}}/4} {
  \node[box,fill=omThm!10,minimum width=2.5cm] (b\c) at (\x,-0.4) {\l};
  \draw[<->] (b\c.north) -- (b\c.north|-e.south);}
\node[font=\scriptsize\itshape] at (2.4,-1.1) {adapters: one per level, backtesters unchanged};
\node[box,fill=gray!10,minimum width=4.4cm] (d) at (0.2,-2.0) {data-access layer\\ named datasets, content hashes};
\node[box,fill=gray!10,minimum width=4.4cm] (r) at (5.4,-2.0) {results store\\ manifest, fills, events per run};
\draw[->] (d.west) -| ([xshift=-4mm]e.west) -- (e.west);
\draw[->] (s.south) -- (s.south|-e.north);
\draw[->] (e.east) -| ++(0.5,-3.2) -- (r.east);
\end{tikzpicture}

\omcaption{The engine's architecture. One strategy object runs at any level; each level's adapter translates between the engine's
context and one of the firm's backtesters, which are wrapped and not edited. Data come through the \omterm{def:pl:backtest-engine-architecture:dal}{data-access layer} with a content hash,
and every run lands in the \omterm{def:pl:backtest-engine-architecture:results}{results store} with its manifest.}\label{fig:pl:backtest-engine-architecture:arch}
\end{omfigure}

The strategy interface has two forms, because the levels differ in what a strategy is. At level~1 there are no orders: a strategy is a
rule that maps closes to a target position, run over all dates at once. At levels~2 to~4 a strategy reacts to the market and to its own
fills by sending and cancelling orders. A \texttt{TargetStrategy} states its rule once, \texttt{target(close, i)}, and the engine runs it
vectorised at level~1 and bar by bar at level~2; a \texttt{QuoteStrategy} has two callbacks, \texttt{on\_market} and \texttt{on\_fill},
and runs at levels~2 to~4. Asking for a quoting strategy at level~1 is an error, not a silently empty backtest.

\omcode{code/platforms/11-backtest-engine-architecture/python/pl_btengine.py}{28}{51}{The chapter's quoting strategy, written once
against the engine's context: one lot at the best bid and ask, requoted when the touch moves, with an inventory limit of three
lots.}\label{lst:pl:backtest-engine-architecture:quoter}

\omcode{code/firm/btengine/firm_btengine.py}{254}{278}{The engine: the level decides the adapter, and every run returns the same
\texttt{Run} with its manifest.}\label{lst:pl:backtest-engine-architecture:engine}

\section{The event model and the clock}

\begin{definition}[Event model]\label{def:pl:backtest-engine-architecture:event}
An \emph{event model}\index{event model} is the set of event classes a \omterm{def:pl:backtest-engine-architecture:engine}{backtest engine} records and the total order in which it
processes and stores them: here each event carries a time, a class and a sequence number, and events are ordered by time, then by a
fixed rank of their class, then by sequence number.
\end{definition}

Four classes suffice for the chapter's strategies, in this rank order: \emph{market} (the top of the book the strategy saw), \emph{fill}
(what the market did to its orders), \emph{order} and \emph{cancel} (what the strategy did in response). The rank settles the question
every event-driven engine must answer and most answer by accident: when a quote update and a fill carry the same time, which does the
strategy see first? A fixed answer makes two runs of the same inputs identical and two levels comparable; the sequence number breaks
the remaining ties in arrival order.

\omcode{code/firm/btengine/firm_btengine.py}{55}{78}{The event model: a total order on (time, class rank, sequence), and the log every
level writes.}\label{lst:pl:backtest-engine-architecture:events}

The clock is the simulated clock of Book~7 (chapter~17): time moves from event to event, never from the machine. Each backtester keeps
its own clock internally --- bar ends at level~2, message times at level~3, nanoseconds of the simulated day at level~4 --- and the
adapters convert all of them to seconds from the start of the session, so that events of different levels can be laid side by side.

\section{Data access}

\begin{definition}[Data-access layer]\label{def:pl:backtest-engine-architecture:dal}
A \emph{data-access layer}\index{data-access layer} is the component through which every backtest reads its data: it resolves a dataset
name to stored data and returns, with the data, an identity of their content (a hash) that the run records.
\end{definition}

The chapter's layer stores named datasets as Parquet (chapters~3 and~4) and keeps a content hash in each file's metadata: a
\texttt{firm.tape} session of Book~7 (the order-by-order messages of a synthetic hour), and daily bars. The hash is what makes ``the same
data'' checkable: two runs that read datasets with equal hashes read the same bytes, whatever the files were called, and a run's
manifest names the hash, not the path. In the firm the same interface fronts the \omterm{def:pl:a-tick-store-on-open-formats:store}{tick store} of chapter~4 and the reference data of
chapter~6; a strategy never opens a file.

\section{The results store}

\begin{definition}[Results store]\label{def:pl:backtest-engine-architecture:results}
A \emph{results store}\index{results store} keeps every backtest run with its manifest --- level, strategy and engine code hashes, data
hashes, parameters and seed --- and its outputs (fills, events, P\&L), under an identifier derived from the manifest, and answers
queries over runs.
\end{definition}

A run's identifier is the hash of its manifest (\cref{lst:pl:backtest-engine-architecture:save}): the same strategy code, engine code,
data, parameters and seed give the same identifier, so a rerun is recognised rather than stored twice, and a changed line of the
strategy gives a new one. The store answers the questions the week of the chapter's hook was spent on: which runs of this strategy
exist, at which levels, on which data, with which code; and, for two runs, where their events first differ.

\omcode{code/firm/btengine/firm_btengine.py}{429}{439}{Saving a run: its manifest, fills and events under the manifest's
hash.}\label{lst:pl:backtest-engine-architecture:save}

\begin{dated}{2026-09}{Open-source engines with one strategy for several environments}\label{dat:pl:backtest-engine-architecture:oss}
Open-source trading platforms describe the same design goal. NautilusTrader's documentation calls it ``an open-source,
production-grade, Rust-native engine for multi-asset, multi-venue trading systems'' in which ``the same strategy and execution-algorithm
code can run across backtest and live systems'', with a common core used by its backtest, sandbox and live contexts and a message bus
between components.
\end{dated}

\section{Adding a level: the reactive simulation}

\begin{definition}[Reactive simulation]\label{def:pl:backtest-engine-architecture:reactive}
A \emph{reactive simulation}\index{reactive simulation} is a backtest in which the strategy's orders are real orders in a simulated
market: they rest in its queues, trade against the other participants' orders and change what those orders execute against afterwards,
so that the market the strategy meets depends on the strategy's presence. The firm uses it as fidelity level~4, the exchange simulator
standing in for the venue.
\end{definition}

At level~3 the strategy's orders are shadows (Book~7, chapter~18): they take a place in the replayed queue and fill when the recorded
executions reach them, but the replay does not change --- a market order that filled a shadow also fills, in the recording, the orders
it filled that day. At level~4 the same \texttt{firm.tape} session is replayed through Book~10's simulator as real orders: each recorded
addition becomes a limit order, each cancellation a cancel, each trade a market order of its size. The strategy's orders are among them.
A market order that meets the strategy's order trades with it and not with the order behind it, which then rests longer and trades later
with something else; when the orders at the touch leave, the strategy's order is left alone at the front, visible to everyone, where the
replay's shadow would have vanished with the level it was hiding in.

\omcode{code/firm/btengine/firm_btengine.py}{378}{410}{Level~4: the dataset's session as the simulator's background flow, the strategy
behind Book~10's \texttt{ReplayStrategyAdapter}, fills and events recorded as at every other
level.}\label{lst:pl:backtest-engine-architecture:level4}

Level~4 is the counterfactual impact of Book~7 (chapter~18) made concrete, with its limits: the background flow is a recording, so its
limit orders do not requote around the strategy and its informed traders do not change their minds. It measures the mechanical part of
the strategy's presence --- queues and matching --- and not the strategic part, which needs agents (Book~10's \texttt{firm.agentmkt}).

\section{The P\&L ladder}

The chapter's experiment runs the quoting strategy of \cref{lst:pl:backtest-engine-architecture:quoter} on twenty one-hour sessions of
\texttt{firm.tape} (seeds 1 to 20) at levels~2, 3 and~4, with 20~$\mu$s of data and order-entry latency and no fees at any level. At level~2
the strategy sees the touch at the end of each one-second bar that has trades and a quote fills in full when the bar trades at its price
(the touch fill of Book~7, chapter~17); at level~3 it takes its place at the back of the queue; at level~4 it is in the queue.

\begin{omfigure}
\begin{tikzpicture}
\begin{axis}[width=11cm,height=5.6cm,ybar,bar width=26pt,ylabel={P\&L per hour (\$)},ymin=-150,ymax=300,ytick={-100,0,100,200,300},
  xtick={2,3,4},xticklabels={{level 2: bars},{level 3: replay},{level 4: reactive}},enlarge x limits=0.3,
  tick label style={font=\small},label style={font=\small},grid=major,grid style={gray!20},table/col sep=comma,
  extra y ticks={0},extra y tick style={grid style={black,thin}}]
\addplot[fill=omDef!45,draw=omDef,error bars/.cd,y dir=both,y explicit]
  table[x=level,y=pnl,y error=se]{figdata/platforms/11-backtest-engine-architecture/ladder.csv};
\end{axis}
\end{tikzpicture}

\omcaption{The quoting strategy's P\&L per hour at three fidelity levels, mean over twenty one-hour sessions with one standard error.
Touch fills on bars make it rich, the replay's queue makes it flat, and a market in which its orders are real makes it lose. Data:
\texttt{fig\_btengine.py}.}\label{fig:pl:backtest-engine-architecture:ladder}
\end{omfigure}

\Cref{fig:pl:backtest-engine-architecture:ladder} is the ladder. At level~2 the strategy makes \$228 an hour ($\pm 29$, one standard
error over the sessions) from 841 fills: a quote that fills whenever the bar trades at its price never waits in a queue. At level~3 it
loses \$12 an hour ($\pm 13$, indistinguishable from zero) from 333 fills: behind the queue, it fills less and mostly when the price is
about to move through it. At level~4 it loses \$102 an hour ($\pm 18$) from 434 fills. The last step is not noise: level~4 is below
level~3 in every one of the twenty sessions (\cref{fig:pl:backtest-engine-architecture:sessions}), by \$90 an hour on average ($\pm 7.4$).

\begin{omfigure}
\begin{tikzpicture}
\begin{axis}[width=8.4cm,height=7cm,xlabel={level 3 (replay): P\&L per hour (\$)},ylabel={level 4 (reactive): P\&L per hour (\$)},
  xmin=-260,xmax=110,ymin=-260,ymax=110,grid=major,grid style={gray!20},tick label style={font=\small},label style={font=\small},
  table/col sep=comma,legend style={font=\footnotesize,at={(0.03,0.97)},anchor=north west}]
\addplot[only marks,mark=*,mark size=1.8pt,omDef] table[x=pnl3,y=pnl4]{figdata/platforms/11-backtest-engine-architecture/sessions.csv};
\addlegendentry{one session}
\addplot[black,dashed,domain=-260:110] {x};
\addlegendentry{equal P\&L}
\end{axis}
\end{tikzpicture}

\omcaption{Each of the twenty sessions at level~3 and level~4. Every point lies below the diagonal: the same strategy on the same
recorded hour does worse when its orders are part of the market. Data: \texttt{fig\_btengine.py}.}\label{fig:pl:backtest-engine-architecture:sessions}
\end{omfigure}

Because the strategy object is one and the same at every level, the difference between two runs is the market's; the engine's event
logs say where it begins. \texttt{first\_divergence} compares the fills of the level-3 and level-4 runs of each session: they part within
the first five fills in all twenty sessions, at a median of 11~seconds into the hour. From then on the two runs are different histories of
the same session.

\section{Where the gap comes from}

Which fills make the difference? Each level-4 fill either has a level-3 counterpart --- a fill of the same side within a millisecond ---
or it exists only because the strategy's order was in the book. Symmetrically, some replay fills have no reactive counterpart. The
ten-second markout of each fill (the mid ten seconds later against the fill price, signed) measures what the fill earned before the
inventory it created was marked; \cref{fig:pl:backtest-engine-architecture:decomposition} splits the markouts of both runs into the three
groups, in the manner of the divergence decomposition of Book~7 (chapter~19).

\begin{omfigure}
\begin{tikzpicture}
\begin{axis}[width=11cm,height=5.6cm,ybar,bar width=16pt,ylabel={10-s markout per hour (\$)},ymin=-80,ymax=30,ytick={-75,-50,-25,0,25},
  xtick=data,symbolic x coords={shared fills,replay-only fills,reactive-only fills},enlarge x limits=0.25,
  tick label style={font=\small},label style={font=\small},grid=major,grid style={gray!20},table/col sep=comma,
  legend style={font=\footnotesize,at={(0.03,0.05)},anchor=south west},extra y ticks={0},extra y tick style={grid style={black,thin}}]
\addplot[fill=omMeth!45,draw=omMeth] table[x=part,y=replay]{figdata/platforms/11-backtest-engine-architecture/decomposition.csv};
\addlegendentry{level 3 (replay)}
\addplot[fill=omThm!45,draw=omThm] table[x=part,y=reactive]{figdata/platforms/11-backtest-engine-architecture/decomposition.csv};
\addlegendentry{level 4 (reactive)}
\end{axis}
\end{tikzpicture}

\omcaption{Ten-second markouts of the two runs' fills, per hour, averaged over the twenty sessions, split into fills both runs share,
fills only the replay has, and fills that exist only at level~4. The shared fills earn the same in both; the replay-only fills earn nothing;
the reactive-only fills lose. Data: \texttt{fig\_btengine.py}.}\label{fig:pl:backtest-engine-architecture:decomposition}
\end{omfigure}

\begin{example}[Four backtesters, one strategy]\label{ex:pl:backtest-engine-architecture:gap}
Per session, 193 fills are shared by the two runs; they earn \$17.5 of ten-second markout at level~3 and \$18.3 at level~4, the same
within noise --- the check that the classification pairs like with like. The replay has 140 fills of its own, which earn \$0.03; the
reactive run has 241 of its own, which lose \$69.2. The markout gap between the runs is \$68.4 an hour, and the reactive-only fills account
for 101\% of it: the loss is made of trades that happen only because the strategy's orders were there: an order that absorbs an aggressive order the
recording gave to someone else, or that is still at the touch, visible, after the orders around it have left. The remaining \$22 of the \$90
gap in P\&L is the inventory those fills leave, marked after more than ten seconds.
\end{example}

\begin{remark}[What the engine settled]\label{rem:pl:backtest-engine-architecture:settled}
With four backtesters and four ports of the strategy, the gap between two levels had three candidate causes: the market, the port and the
backtester's conventions (what the strategy sees, when a fill is reported). The engine removes the second by construction and makes the
third explicit in the adapters. What remains is what the levels are for.
\end{remark}

\section{A rule at levels 1 and 2}

The same engine runs a daily rule --- long when the close is above its level fifty days earlier, short when below --- on a thousand days of
synthetic bars at level~1 (\texttt{firm.vecbt}, the rule's weight earned from the next close) and at level~2 (\texttt{firm.evbt}, market
orders filled at the next open). Level~1 returns $-26.5\%$ over the period and level~2 $-31.5\%$; the five points are the overnight moves
on the days the position changes, earned by the old position at level~2 and by the new one at level~1, which trades at a close it has
only just seen. The rule is written once; the two runs sit next to each other in the \omterm{def:pl:backtest-engine-architecture:results}{results store} under different identifiers.

\section{Tutorial: one strategy through four backtesters}\label{tut:pl:backtest-engine-architecture:tut}

\begin{tutorial}
\textbf{Goal.} Run one strategy object at every level through one engine, store every run, and find where two levels part.
\textbf{End state:} \cref{fig:pl:backtest-engine-architecture:ladder,fig:pl:backtest-engine-architecture:decomposition} and a \omterm{def:pl:backtest-engine-architecture:results}{results
store} holding sixty runs.
\begin{tutsteps}
\item \textbf{Data}: \texttt{DataAccess.tape(3600, seed)} stores a session with its hash; \texttt{daily()} stores the daily bars.
\item \textbf{One session}: \texttt{pl\_btengine.one\_session} runs levels~2, 3 and~4 and saves each run.
\item \textbf{The ladder}: \texttt{ladder()} over seeds 1 to 20, \texttt{summary(rows)}.
\item \textbf{The gap}: \texttt{classify\_fills}, \texttt{markouts} and \texttt{first\_divergence} on the level-3 and level-4 runs.
\item \textbf{Identity}: \texttt{rerun\_is\_identical()}: same inputs, same run identifier, same events.
\end{tutsteps}
\textbf{What to change next.} Add a one-tick spread condition to the strategy (quote only when the spread is one tick) and see which
level's verdict changes; run level~4 with Book~10's agent population instead of the recorded session.
\end{tutorial}

\section{Build: the backtest engine}\label{bld:pl:backtest-engine-architecture:btengine}

\begin{build}
\bfield{Purpose} One engine over the firm's four fidelity levels, so that a strategy is written once, runs anywhere on the ladder, and
every run can be found, compared and reproduced. Chapter~12 hosts the same strategy objects in production-shaped environments.
\bfield{Interface} \texttt{Engine(level, strategy, data, config).run() -> Run}; \texttt{TargetStrategy}, \texttt{QuoteStrategy};
\texttt{Event}, \texttt{EventLog}; \texttt{DataAccess(root)} with \texttt{put}, \texttt{get}, \texttt{hash}, \texttt{tape},
\texttt{daily}; \texttt{ResultsStore(root)} with \texttt{save}, \texttt{load}, \texttt{query}; \texttt{first\_divergence},
\texttt{classify\_fills}, \texttt{markouts}.
\bfield{Rules} Book~7's backtesters and Book~10's simulator are wrapped, never edited; events in the total order (time, class rank,
sequence); a run's identifier is its manifest's hash; a strategy never opens a file; no level silently accepts a strategy it cannot run.
\bfield{Acceptance tests} \texttt{code/firm/btengine/tests/}: the event order on simultaneous events; content hashes; level~1 equal to
\texttt{firm.vecbt} called directly, level~3 to \texttt{firm.lobreplay}, level~4 to \texttt{firm.exchsim} with the adapter; equal run
identifiers for equal inputs and different ones for a changed parameter; queries; the fill classification and markouts on a worked case.
\bfield{Stretch} Parallel runs over seeds (chapter~13); an agent-population background at level~4; a level-3 run from the \omterm{def:pl:a-tick-store-on-open-formats:store}{tick store}'s
recorded feed of chapter~4.
\end{build}

\begin{omsources}
\item One Quant Book~7, chapters 16--19 (fidelity levels, event-driven backtests, order-book replay, simulation against live); One Quant
Book~10, chapter~26 (the exchange simulator).
\item NautilusTrader documentation, ``Overview'' and ``Architecture''.
\end{omsources}

\section{Exercises}

\begin{exercise}[$\star$]\label{exo:pl:backtest-engine-architecture:1}
Order these simultaneous events by the engine's \omterm{def:pl:backtest-engine-architecture:event}{event model}: an order sent at $t = 3$, a quote update at $t = 3$, a fill at $t = 3$, a
cancel at $t = 2$.
\end{exercise}

\begin{exercise}[$\star$]\label{exo:pl:backtest-engine-architecture:2}
Why does the engine refuse a quoting strategy at level~1 instead of running it with no fills?
\end{exercise}

\begin{exercise}[$\star$]\label{exo:pl:backtest-engine-architecture:3}
Two runs have the same strategy code, data hashes and seed but different run identifiers. What can differ?
\end{exercise}

\begin{exercise}[$\star\star$]\label{exo:pl:backtest-engine-architecture:4}
From the ladder, is the level-3 loss of \$12 an hour ($\pm 13$) evidence that the strategy loses in replay? Is the level-3 to level-4 gap
evidence of anything?
\end{exercise}

\begin{exercise}[$\star\star$]\label{exo:pl:backtest-engine-architecture:5}
Explain why the shared fills should earn the same markout at levels~3 and~4, and what it would mean if they did not.
\end{exercise}

\begin{exercise}[$\star\star$]\label{exo:pl:backtest-engine-architecture:6}
What does level~4 as built here not capture about the strategy's effect on the market, and which component would add it?
\end{exercise}

\begin{exercise}[$\star\star\star$]\label{exo:pl:backtest-engine-architecture:7}
\emph{Coding.} Run \texttt{one\_session} for seed~3 with the inventory limit raised from three lots to ten. How do the three levels'
P\&L change, and does the reactive-only share of the gap move?
\end{exercise}

\begin{exercise}[$\star\star\star$]\label{exo:pl:backtest-engine-architecture:8}
\emph{Find the flaw.} ``We compared the strategy in our replay backtester and in the new simulator and the simulator was worse, so the
simulator's matching has a bug. We ported the strategy carefully to both.''
\end{exercise}

\section{Problem: Four Backtesters, One Strategy}

\begin{problem}[{Weekend problem --- the ladder and its gap}]\label{pb:pl:backtest-engine-architecture:1}
The quoting strategy of \cref{lst:pl:backtest-engine-architecture:quoter}, twenty one-hour sessions, levels~2 to~4.

\medskip\noindent\textbf{Part I --- The engine.}
\begin{enumerate}
\item What must a \omterm{def:pl:backtest-engine-architecture:engine}{backtest engine} own that each backtester does not?
\item Why two strategy forms, and which levels run each?
\item What does the class rank of the \omterm{def:pl:backtest-engine-architecture:event}{event model} settle?
\item What does a run's manifest contain, and why is the identifier its hash?
\item Why does the \omterm{def:pl:backtest-engine-architecture:dal}{data-access layer} return a hash with the data?
\end{enumerate}

\medskip\noindent\textbf{Part II --- The levels.}
\begin{enumerate}[resume]
\item What does the strategy see, and when does a quote fill, at level~2?
\item How does a level-3 shadow order differ from a level-4 order?
\item What does level~4 replay from the session, and as what?
\item What does level~4 not capture?
\item How is Book~7's level~4 represented here?
\end{enumerate}

\medskip\noindent\textbf{Part III --- The ladder.}
\begin{enumerate}[resume]
\item What are the P\&L per hour and the fills at each level?
\item Which differences are statistically clear?
\item In how many sessions is level~4 below level~3, and by how much on average?
\item When do the level-3 and level-4 runs first differ?
\item Why are the level-2 fills so many?
\end{enumerate}

\medskip\noindent\textbf{Part IV --- The gap.}
\begin{enumerate}[resume]
\item State the \emph{named result}: the P\&L an hour at each level and the share of the replay-to-reactive gap explained by fills that
exist only because the strategy was in the book.
\item Which fills lose, and why?
\item What part of the gap is not in the ten-second markouts?
\item What would you conclude about the strategy?
\item In one sentence: what does the engine make possible that four backtesters did not?
\end{enumerate}
\end{problem}

\section{Interview questions}

\begin{interviewq}[$\star$ \iqroles{researcher}]\label{iq:pl:backtest-engine-architecture:1}
Your strategy is profitable in the bar backtest and loses in the order-book simulation. Where do you look?
\end{interviewq}

\begin{interviewq}[$\star\star$ \iqroles{developer}]\label{iq:pl:backtest-engine-architecture:2}
Two events in your event-driven backtester have the same timestamp. How should the engine order them?
\end{interviewq}

\begin{interviewq}[$\star\star$ \iqroles{developer, researcher}]\label{iq:pl:backtest-engine-architecture:3}
How would you make every backtest run reproducible and findable a year later?
\end{interviewq}

\begin{interviewq}[$\star\star$ \iqroles{researcher}]\label{iq:pl:backtest-engine-architecture:4}
What does a replay of the order book miss for a passive strategy, and how would you measure it?
\end{interviewq}

\begin{interviewq}[$\star\star\star$ \iqroles{developer}]\label{iq:pl:backtest-engine-architecture:5}
Design a \omterm{def:pl:backtest-engine-architecture:engine}{backtest engine} that covers vectorised backtests, event-driven backtests and a market simulator with one strategy API.
\end{interviewq}
